\documentclass[11pt,a4paper]{article}
\usepackage[utf8]{inputenc}
\usepackage[T1]{fontenc}
\usepackage[spanish]{babel}
\usepackage{hyperref}
\usepackage{geometry}
\geometry{margin=2.5cm}
\usepackage{parskip}

\title{009 el cierre}
\author{Alejandro Valenzuela Rivera\\\small Universidad Autónoma de Nuevo León}
\date{\today}
\begin{document}
\maketitle

\begin{quote}
\textbf{Sesgo declarado.} Este trabajo se sostiene desde el sesgo declarado de su autor, y no finge neutralidad. Lo que afirma se puede tumbar, y ese es el punto: el falsificador está escrito, no escondido.
\end{quote}

\section*{Resumen}
La deuda se paga tarde; no se cierra con arrepentimiento sino con admision.

\section{Por qué se sostiene}
El arrepentimiento deja limpio al autor; la admision carga al otro.

\section{De dónde viene}
Breaking Bad.

\section{El caso}
Breaking Bad, Walter White: admitir en vez de arrepentirse.

\section{Método}
Un trabajo que cierra con lo que quedo abierto declarado.

\section{Límites, y qué lo tumba}
Un trabajo que se cierra sin decir que quedo sin hacer.

\section{Lo que queda abierto}
\textquestiondown{}Por quien lo estoy haciendo?

\section*{Cómo se leyó esta tesis}
Esto se escribió en el repositorio man-made-empathy, del que este candado es una pieza: candados/06 del caso/009 el cierre.md. La máquina que lo compuso lo hizo desde este mismo archivo, y no agregó nada que no estuviera escrito. Si algo de esto es falso, se dice en el mismo sitio, y se corrige con fecha.

\end{document}
